\chapter{Wstęp}
\label{chapt:Wstep}

Tutaj napisz krótkie wprowadzenie do swojej pracy. Możesz omówić podstawowe algorytmy związane z celem swojej pracy. Możesz podać krótki rys historyczny. Pamiętaj jednak o zwięzłości tekstu i o unikaniu nic nie wnoszących ogólników.   
Oto przykładowy początek takiego wprowadzenia:

\begin{quote}
Problem algorytmów wyspecjalizowanych do analizy danych strumieniowych wzbudził duże zainteresowanie badaczy na początku XX wieku. Jednen z pierwszych takich algorytmów służył do wyboru losowego elementu zgodnie z rozkładem jednostajnym  z okna złożonego z ostatnich $N$ elementów strumienia, gdzie $N$ jest ustaloną liczbą.  ...
\end{quote}

      
\section{Cel i struktura pracy}
                                       
Postaraj się tutaj opisać możliwie zwięźle i precyzyjnie cel pracy. Oto przykład początku tej sekcji pracy:

\begin{quote}
W pracy tej będziemy zajmować się problemem wyznaczania $k$ - najczęściej występujących elementów w oknie złożonym z $N$ ostatnich elementów strumienia danych. ... 
\end{quote}

\subsection*{Struktura pracy}

Tutaj opisz krótko, po kolei, wszystkie rozdziały swojej pracy. Oto przykład:

\begin{quote}
Rozdział pierwszy zaczynamy (patrz \ref{chapt1:sec:intro}) od opisu wymagań stawianych algorytmom strumieniowym (główne wymaganie to wymaganie użycia logarytmicznej pamięci względem rozmiaru strumienia), opisujemy pojęcie przesuwnego okna strumienia danych, omawiamy algorytm Bravermana,  ... .\par

W rozdziale drugim omawiamy ... \par
...
\end{quote}

\subsection*{Zastosowane technologie}

Tutaj opisz, z jakich języków programowania oraz z jakich kompilatorów korzystałeś, 
w jakim środowisku uruchamiałeś oraz testowałeś budowane oprogramowanie. Wskaż również najważniejsze biblioteki wykorzystane w zbudowanych programach.

\section{Wykorzystanie AI}

Jeśli podczas realizacji pracy wykorzystywano współczesne narzędzia AI, 
to w tej części należy je opisać, czyli napisać, jakie narzędzia AI były wykorzystywane, do czego były używane oraz których miejsc pracy to dotyczy. \textbf{Tej części pracy nie możesz pominąć}, aczkolwiek możesz ją umieścić w innym miejscu pracy. 
Koniecznie zapoznaj się z ``Wytycznymi dotyczącymi zasad wykorzystywania AI'', które znajdują się na stronach naszego Wydziału (patrz \cite{WITWytyczneSI}).\par

Jeśli AI było wykorzystywane do generowania kodów, to również o tym tutaj napisz, a w kodach programu (w komentarzach) zaznacz  te fragmenty, w które ingerowała AI. Oto przykład:

\begin{quote}
  W pracy korzystano z asystenta Gemini 3.0 firmy Google w wersji Pro do pomocy w przekształceniu kodów z języka programowania C na język programowania Haskell (patrz, między innymi, \cref{lst:hhaskell:qs}) oraz do pomocy w dokumentowaniu kodów wszystkich zbudowanych w trakcie realizacji pracy modułów. 
\end{quote}

Jeśli podczas realizacji pracy nie była  wykorzystywana AI, to w sekcji tej umieść mniej więcej takie zdanie: 
"\textit{W pracy nie korzystano ze Sztucznej Inteligencji}".    


\section{Wprowadzenie} 
\label{chapt1:sec:intro}

Tutaj możesz opisać podstawowe pojęcia, terminy i oznaczenia, którymi posługujesz się w pracy. 

\section{Podstawy teoretyczne}

Możesz podać podstawowe wyniki, z których korzystasz w pracy. 
Mogą to być odwołania do literatury - w rozdziale \ref{examples:cite} omówiono sposób formatowania bibliografii oraz metody odwoływania się w pracy do konkretnych pozycji (\verb|\cite{...}|, \verb|\citeauthor{...}|, ...)    

